\documentclass[11pt]{article}
\usepackage[a4paper,margin=25mm,headheight=14pt]{geometry}
\usepackage[T1]{fontenc}
\usepackage[utf8]{inputenc}
\usepackage{lmodern,microtype}
\usepackage{amsmath,amssymb,amsthm,mathtools}
\usepackage{booktabs,longtable,tabularx,array}
\usepackage{enumitem}
\usepackage{xcolor}
\usepackage{listings}
\usepackage{fancyhdr}
\usepackage[hidelinks]{hyperref}
\hypersetup{pdftitle={Bounded-interface integral homology for compressed unknot hierarchies},pdfauthor={Research draft prepared for ProveIt with ChatGPT}}
\pagestyle{fancy}\fancyhf{}
\fancyhead[L]{\small Bounded-interface integral homology}
\fancyhead[R]{\small ProveIt research continuation}
\fancyfoot[C]{\thepage}
\setlength{\parskip}{4pt}
\setlength{\parindent}{0pt}
\setlist{itemsep=3pt,topsep=4pt,leftmargin=*}
\lstset{basicstyle=\ttfamily\small,breaklines=true,columns=fullflexible,keepspaces=true,frame=single,framerule=.3pt}
\newtheorem{theorem}{Theorem}[section]
\newtheorem{lemma}[theorem]{Lemma}
\newtheorem{proposition}[theorem]{Proposition}
\newtheorem{corollary}[theorem]{Corollary}
\theoremstyle{definition}
\newtheorem{definition}[theorem]{Definition}
\newtheorem{example}[theorem]{Example}
\theoremstyle{remark}
\newtheorem{remark}[theorem]{Remark}
\newcommand{\Z}{\mathbb Z}
\newcommand{\Q}{\mathbb Q}
\newcommand{\F}{\mathbb F_2}
\newcommand{\coker}{\operatorname{coker}}
\newcommand{\Tor}{\operatorname{Tor}}
\newcommand{\rank}{\operatorname{rank}}
\newcommand{\im}{\operatorname{im}}
\newcommand{\code}[1]{\texttt{#1}}
\newcommand{\satO}{\mathcal O}
\newcommand{\satN}{\mathcal N}
\newcommand{\bits}{\operatorname{bits}}
\title{Bounded-interface integral homology\\[5pt]
{\large Binary surface condensation, twisted attachments, and\\
provenance-aware terminals for compressed unknot hierarchies}}
\author{Research draft prepared for the ProveIt project\\\small Prepared with ChatGPT}
\date{12 October 2026}
\begin{document}
\maketitle

\begin{abstract}
The current ProveIt implementation already constructs a normal cut in binary
cell families and extracts its all-dimensional source parallelity bundle,
leaving a bounded inherited handle core. We develop the next algebraic layer:
exact integral homology from that core and its bounded interface, without
expanding the genera, boundary counts, or multiplicities of the attached
surface bundles. The interacting part of $H_1$ is the cokernel of a matrix
whose dimensions depend only on the core and the number of attached boundary
circles. All other free ranks and repeated $2$-torsion are binary counters.
A saturated nonorientable base contributes one generator and one relation
$2h+s=0$, not one generator per crosscap. This yields a sharp mod-two torsion
formula and a factor-two bound on torsion-exponent growth.

We also give bounded marked-port reconstruction, a torsion bound in terms
of actual interface incidence, and compact cohomology witnesses. For genuine
cut descendants of a certified knot exterior in $S^3$, a classical duality
argument bypasses interior homology entirely: boundary genera determine all
integral homology groups. Combined with the condensation theorem, this gives
primitive-lattice and mod-two independence conditions on the attachments.
A single sphere boundary is a ball terminal only under the embedding
provenance, not for a generic torus-boundary input.

The accompanying standard-library implementation has independent algebraic
replay, a checked small Smith kernel, and 4,800 seeded comparisons across six
validation families. Paired complete homology-kernel measurements record
both large-input gains and small-input regressions. These are interface-model
measurements, not knot-recognition benchmarks. The source-geometric adapter,
constructive cocycle lift, admissible patterned hierarchy, and global
quasi-polynomial bound are not supplied by this contribution.
\end{abstract}

\tableofcontents
\newpage

\section{Research position and the concrete bottleneck}
\label{sec:position}

\subsection{The audited starting point}
This continuation is based on ProveIt revision
\begin{center}\small\code{ed166080d7db14dd7a91d07788996f02947dd92a}.\end{center}
The revision was committed on 12 October 2026. The relevant implementation is
\path{Topology/UnknotRecognition/fast/}; the maintained theoretical account is
under \code{synthesis/}. The source audit in Appendix~\ref{app:audit} records
specific paths and Git blob identifiers rather than treating the entire
repository as a single reviewed object.

Two prior achievements are essential. The raw-cut constructor preserves all
cell dimensions and source attaching charts in affine binary families
\cite{repo-cells}. More recently, \code{normal\_cut\_parallelity} extracts
product handles in indices zero, one, and two, including cases where a base
component has no selected zero- or one-handles. It retains base surface types,
fibre twisting, and core-facing ports. Its inherited complementary core has
at most $30T$ handles, with a conservative bound of $66T$ core ports, for $T$
source tetrahedra \cite{repo-parallelity}. These results are existing work,
not contributions of the present article.

The inherited core is not yet an admissible patterned hierarchy structure.
Likewise, a binary polygon description is not automatically a polynomial-time
homology algorithm: an explicit chain matrix can have a number of rows
proportional to the numerical normal weight. Our question is more specific:
\emph{which part of integral homology actually needs an explicit matrix once
the core--bundle interface has been supplied?}

\subsection{Relation to the published hierarchy algorithms}
Polynomial-time extraction of parallelity data is not an absent theorem.
Lackenby's Theorem 9.3 in \emph{The efficient certification of knottedness and
Thurston norm} computes the complementary core and the types and vertical
attachments of parallelity components for an orientable normal surface, in
time polynomial in the tetrahedron count and logarithmic normal weight
\cite{lackenby-cert}. That theorem and the current repository's all-dimensional
source implementation are the proper antecedents of this work.

Lackenby's July 2026 manuscript gives an iteration bound $L(g+1)^L$, where
$L$ bounds hierarchy length and $g$ bounds the relevant pattern complexity.
It separately identifies the need to specify and bound the cost of individual
operations \cite[Section 9]{lackenby-2026}. We address one such cost, not the
entire search. The local algebra below cannot by itself bound the number of
hierarchy corrections, the size of every later patterned core, or the cost
of discovering the necessary surfaces.

\subsection{What is new here, and what is classical}
Mayer--Vietoris, the homology presentations of compact surfaces, Smith normal
form, Alexander duality, and the Schoenflies theorem are classical
\cite{hatcher,brown,kannan-bachem}. The mathematical contribution here is their
explicit bounded-interface organization: a signed integral matrix, a complete
binary bulk formula, one additional generator per saturated twisted base,
and interface conditions that are directly checkable downstream of source
replay. We prove the resulting statements in full and implement the algebraic
contracts. We do not claim that these consequences are previously unknown in
every formulation in the literature.

The most immediately usable outcomes are the reduction in
Theorem~\ref{thm:condensation}, the bounded-port extraction of
Theorem~\ref{thm:ports}, the twisted-attachment theorem
\ref{thm:twisting}, and the source-aware dispatch rule in
Section~\ref{sec:s3}. The proofs remain useful even when a particular
implementation chooses another normal-form backend.

\section{The geometric and algebraic input contracts}
\label{sec:contract}

\subsection{A core with interval bundles}
Let $X$ be a compact orientable three-manifold, possibly disconnected, with
nonempty boundary on every component. Assume a decomposition
\begin{equation}
 X=K\cup\bigcup_i E_i,
 \qquad E_i\longrightarrow B_i,
 \qquad K\cap E_i=E_i|_{A_i},
 \label{eq:decomp}
\end{equation}
where $K$ is a finite core, each $E_i$ is an interval bundle over a compact
connected surface $B_i$, and $A_i$ is a compact one-dimensional submanifold
of $\partial B_i$. Distinct bundle components are disjoint. Collars and
corners are arranged so that these intersections give the usual
Mayer--Vietoris sequence. In particular, there is no unrecorded triple
intersection. A source certificate must justify this contract when the
input comes from a normal cut.

Each component of $A_i$ is either a closed interval or an entire boundary
circle of $B_i$. Write $a_i$ and $c_i$ for their respective counts, and
$b_i=|\pi_0(\partial B_i)|$. Adjacent core-facing arcs must first be merged
through their actual quotient vertices. Counting raw port arcs as components
of $A_i$ is generally wrong.

For an orientable base let $g_i$ be its genus. For a nonorientable base let
$k_i\geq1$ be its crosscap number. Thus
\[
 \chi(B_i)=2-2g_i-b_i\quad\hbox{or}\quad2-k_i-b_i.
\]
An orientable total interval bundle over a nonorientable base is twisted.
The homology reduction uses the base's homotopy type; it does not infer
bundle twisting from the slope of a singleton affine map.

\begin{definition}
A base with $b_i>0$ is \emph{saturated} if $c_i=b_i$: its entire boundary is
attached through annuli. A saturated base has $a_i=0$. Write $\satO$ and
$\satN$ for the saturated orientable and nonorientable bases, and put
$o=|\satO|$, $t=|\satN|$.
\end{definition}

Unattached bundle components can be grouped in profiles with a binary
multiplicity $m_i$. An attached component is individually identified and
has $m_i=1$. In particular, a port cannot silently refer to several identical
copies. No loop in the intended computation is indexed by $g_i$, $k_i$,
$b_i$, or $m_i$.

\subsection{The core presentation and signed circle maps}
Take a finite CW model for $K$, with $v$ vertices, $e$ edges, $f$ two-cells,
and $u$ three-cells. A spanning forest in its one-skeleton gives
\[
 H_1(K;\Z)=\Z^d/\im R,\qquad d=e-v+b_0(K),
 \qquad R\in\Z^{d\times f}.
\]
Every core two-cell has a closed edge word. Its column in $R$ is the exponent
sum on the nonforest edges. Collapsing the forest simply deletes tree-edge
coordinates; it does not introduce long word substitutions.

For the attached boundary circles of $B_i$, let
\[
 z_{i1},\ldots,z_{ic_i}\in\Z^d
\]
be their image cycles in this presentation. For an orientable base, orient
all its boundary components using one orientation of the base. This
\emph{relative} orientation information is indispensable. Local consistency
around each individual circle does not establish consistency between different
boundary circles of the same orientable base.

For a nonorientable base, use the boundary convention in the presentation
\begin{equation}
 2(x_1+\cdots+x_{k_i})+d_1+\cdots+d_{b_i}=0.
 \label{eq:nonorientable}
\end{equation}
Changing one boundary direction changes the eventual sum by $2z$ and is
absorbed by a change of the retained generator. For saturated bases define
\begin{equation}
 s_i=\sum_{j=1}^{b_i}z_{ij}.
 \label{eq:boundary-sum}
\end{equation}
A source implementation must provide these cycles, not just their number,
a port mask, or an Euler characteristic.

\subsection{The attachment graph}
Let $\Gamma$ have one vertex for every connected component of $K$ and every
attached base. Each component of $A_i$ contributes one edge joining its base
vertex to the component of $K$ containing its image. Parallel edges are
retained. Then
\begin{equation}
 r_\Gamma=|E(\Gamma)|-|V(\Gamma)|+b_0(\Gamma).
 \label{eq:graph}
\end{equation}
Unattached bundles are not vertices of $\Gamma$; they are counted separately.
The empty graph has all three counts zero. This convention handles an empty
core without a special topological argument.

\subsection{What the delivered input validator establishes}
The Python model checks the integer domains, closed and contiguous core edge
words, attachment counts, multiplicities, and graph incidence. These checks
make a valid algebraic graph-of-spaces input. They do \emph{not} establish
that a native raw-cut packet realizes it, that $K$ has the claimed
three-manifold attaching maps, or that $X$ embeds in $S^3$. Those are separate
geometric premises. The homology formula for the graph of spaces is exact;
its interpretation as the homology of a particular source cut requires
source binding.

\section{Exact integral condensation}
\label{sec:condensation}

\subsection{The small matrix and binary bulk ranks}
List the columns $s_i$ for $i\in\satO$ and $i\in\satN$ as matrices $S_O$
and $S_N$. Set
\begin{equation}
 C=\begin{pmatrix}
 R&S_O&S_N\\
 0&0&2I_t
 \end{pmatrix}\in\Z^{(d+t)\times(f+o+t)}.
 \label{eq:matrix}
\end{equation}
The extra row for a saturated nonorientable base corresponds to a generator
$h_i$, and its column asserts $s_i+2h_i=0$.

Define the per-copy free contribution $\lambda_i$ by Table~\ref{tab:local}.
The cases $b_i=0$ are necessarily unattached. The last column records only
the isolated closed nonorientable contribution, not torsion that might be
present in $\coker C$.

\begin{table}[ht]
\centering
\begin{tabular}{@{}lll@{}}
\toprule
Base and attachment condition & $\lambda_i$ & Extra $\Z/2$ per copy\\
\midrule
Orientable, $b_i>c_i$ & $2g_i+b_i-1-c_i$ & $0$\\
Nonorientable, $b_i>c_i$ & $k_i+b_i-1-c_i$ & $0$\\
Orientable, $b_i=c_i>0$ & $2g_i$ & $0$\\
Nonorientable, $b_i=c_i>0$ & $k_i-1$ & $0$\\
Orientable, $b_i=0$ & $2g_i$ & $0$\\
Nonorientable, $b_i=0$ & $k_i-1$ & $1$\\
\bottomrule
\end{tabular}
\caption{Local surface data retained as binary counters.}
\label{tab:local}
\end{table}

Set
\begin{equation}
 \beta=r_\Gamma+\sum_i m_i\lambda_i,
 \qquad
 \mu=\sum_{\substack{i:\ b_i=0\\ B_i\text{ nonorientable}}}m_i.
 \label{eq:bulk}
\end{equation}
The sum is over the supplied profiles; attached profiles have multiplicity
one. These expressions involve ordinary integer arithmetic, not expanded
lists of basis elements.

\begin{theorem}[Bounded-interface integral homology]
\label{thm:condensation}
Under the contract of Section~\ref{sec:contract},
\begin{align}
 H_0(X;\Z)&\cong\Z^{b_0(\Gamma)+\sum_{A_i=\varnothing}m_i},
 \label{eq:h0}\\
 H_1(X;\Z)&\cong\coker C\ \oplus\ \Z^\beta\ \oplus\ (\Z/2)^\mu.
 \label{eq:h1}
\end{align}
The inclusion of the core is represented by its original $d$ generators in
$\coker C$. No matrix dimension in \eqref{eq:matrix} depends on a numerical
base genus, crosscap number, boundary count, or unattached multiplicity.
The direct-sum splitting in \eqref{eq:h1} is not asserted to be canonical.
\end{theorem}

\begin{proof}
The interval-bundle projection is a homotopy equivalence $E_i\simeq B_i$,
and its restriction is a homotopy equivalence $E_i|_{A_i}\simeq A_i$.
Mayer--Vietoris therefore gives the exact sequence
\begin{multline}
 \bigoplus_i H_1(A_i)\longrightarrow
 H_1(K)\oplus\bigoplus_i H_1(B_i)
 \longrightarrow H_1(X)\\
 \longrightarrow\bigoplus_i H_0(A_i)
 \longrightarrow H_0(K)\oplus\bigoplus_i H_0(B_i).
 \label{eq:mv}
\end{multline}
An interval contributes zero to $H_1(A_i)$; an attached circle identifies its
boundary generator with $z_{ij}$. The kernel of the final map is the integral
cycle group of $\Gamma$, a free group of rank $r_\Gamma$. Thus the middle
cokernel extends by $\Z^{r_\Gamma}$. An extension with free quotient splits
as an abelian group. The same incidence map gives \eqref{eq:h0}.

It remains to compute the middle cokernel explicitly. An orientable genus
$g$ base has the abelian presentation
\[
 \left\langle a_1,b_1,\ldots,a_g,b_g,d_1,\ldots,d_b
 \ \middle|\ d_1+\cdots+d_b=0\right\rangle.
\]
When $b=0$, the displayed relation is empty and its group is $\Z^{2g}$.
For a nonorientable base, use \eqref{eq:nonorientable}. These presentations
follow by abelianizing the standard polygon presentations; the commutator
terms disappear in the orientable case \cite{hatcher}.

Suppose first that $c<b$. Reorder boundary components so that the attached
ones are $d_1,\ldots,d_c$, and $d_b$ is unattached. Eliminate $d_b$ using the
single surface relation. This is an integral elimination with coefficient
one, including in the nonorientable case. Each attached relation then
eliminates $d_j$ in favor of $z_{ij}$. Since there is still the eliminated
unattached boundary generator available, no relation is imposed on the core.
The remaining free counts are $2g+b-1-c$ and $k+b-1-c$.

If $c=b>0$ and the base is orientable, substitute all $d_j=z_{ij}$.
The only residual relation is $s_i=0$, while all $2g$ handle generators
remain free. If the base is nonorientable, introduce
\[
 h=x_1+\cdots+x_k,\qquad y_j=x_j\quad(2\leq j\leq k).
\]
This is a unimodular change of basis, since
$x_1=h-y_2-\cdots-y_k$. The residual relation is $2h+s_i=0$; the $k-1$
generators $y_j$ remain free. This accounts for exactly one new row and one
new column in \eqref{eq:matrix}, irrespective of $k$.

Finally, a closed nonorientable base has
$H_1\cong\Z^{k-1}\oplus\Z/2$, by the same change of basis. An unattached
bundle is a separate component of $X$, so each such copy contributes one
independent $\Z/2$. An unattached orientable closed base contributes only
$\Z^{2g}$. Applying these eliminations component by component gives
\eqref{eq:matrix}, Table~\ref{tab:local}, and \eqref{eq:h1}. All eliminations
retain the original core generators, proving the inclusion assertion.
\end{proof}

\subsection{Recovering the remaining Betti numbers}
The decomposition gives
\begin{equation}
 \chi(X)=v-e+f-u+\sum_i m_i\chi(B_i)-\sum_i a_i,
 \label{eq:euler}
\end{equation}
since a circle has Euler characteristic zero and an interval has Euler
characteristic one. If $r=\rank_\Q C$, then
\begin{equation}
 b_1(X)=d+t-r+\beta.
 \label{eq:b1}
\end{equation}

\begin{corollary}
\label{cor:h2}
For an orientable three-manifold satisfying the nonempty-boundary hypothesis,
\[
 H_3(X;\Z)=0,\qquad
 H_2(X;\Z)\cong\Z^{\chi(X)-b_0(X)+b_1(X)}.
\]
\end{corollary}
\begin{proof}
The top absolute homology vanishes on every component with boundary.
Poincar\'e--Lefschetz duality identifies $H_2(X;\Z)$ with
$H^1(X,\partial X;\Z)$. The latter is torsion-free: the possible universal
coefficient term $\operatorname{Ext}(H_0(X,\partial X),\Z)$ vanishes, and
$\operatorname{Hom}(H_1(X,\partial X),\Z)$ is free for this finite complex.
Euler characteristic now gives the stated rank.
\end{proof}

This corollary must not be applied merely because an arbitrary algebraic
model has an integer called \code{three\_cells}. The delivered software
reports the Euler value but does not certify the manifold premise or claim
an unconditional $H_2$ computation from that count alone.

\subsection{Two small examples with different behavior}
\begin{example}[A saturated annulus]
Let $H_1(K)=\Z\langle z\rangle$. Attach both boundary circles of an annulus
to the same core component. The attachment graph has two parallel edges and
$r_\Gamma=1$. If the coherently oriented images are $z,-z$, then $C=(0)$
and $H_1(X)=\Z^2$. If they are $z,z$, then $C=(2)$ and
$H_1(X)=\Z\oplus\Z/2$. The scalar data $(\chi,b,c)$ are identical.
\end{example}

\begin{example}[A saturated M\"obius base]
For a M\"obius band with attached boundary image $nz$,
\[
 H_1(X)=\langle z,h\mid nz+2h=0\rangle
 \cong\Z\oplus\Z/\gcd(n,2).
\]
Only the parity of $n$ governs the torsion in this example. Replacing this
base by a nonorientable base of $k$ crosscaps adds $k-1$ free generators,
not a larger interacting matrix.
\end{example}

\section{Recovering the bounded attachment interface}
\label{sec:ports}

The repository already retains a bounded set of core-facing arcs, with
source ranks and endpoint maps. The missing algebraic distinction is whether
several such arcs form an interval or an entire boundary circle, and what
that circle represents in the core. This section isolates a bounded
postprocessing task rather than expanding all base polygons.

\subsection{A single marked endpoint query}
Suppose there are $p$ raw core-facing arcs. Give the two endpoint occurrences
of arc $j$ distinct labels $2j$ and $2j+1$, even if their source corner points
coincide. Assign endpoint occurrence $r$ the weight $2^r$. At a source point
carrying several occurrences, add their weights. The quotient-vertex
relation on the full base polygon system is already an interval-orbit
problem in the existing geometry.

\begin{lemma}[Endpoint masks do not alias]
\label{lem:masks}
In the weighted quotient-vertex histogram, the nonzero orbit weights are
exactly the bit masks of the endpoint equivalence classes. Each nonzero
weight occurs on exactly one orbit. There are at most $2p$ nonzero weights,
and every weight has at most $2p$ bits.
\end{lemma}
\begin{proof}
Each endpoint occurrence is marked once. An orbit weight is consequently a
sum of distinct powers of two. Its binary expansion records exactly which
endpoint labels belong to that orbit; no carry or multiplicity ambiguity is
possible. Distinct marked orbits have disjoint nonempty sets of labels, so
cannot have the same nonzero weight. Every marked orbit contains at least
one of the $2p$ endpoints.
\end{proof}

The unmarked orbits can be numerous; their multiplicity is irrelevant here.
The output bound concerns this particular endpoint query, not an arbitrary
weighted histogram with unrelated weights. A replayed orbit calculation is
still needed to prove that the displayed masks are the correct geometric
ones. Disjoint, exhaustive masks by themselves are not such a proof.

\subsection{The finite port graph}
Identify endpoint occurrences according to Lemma~\ref{lem:masks}, and keep
one edge for each raw core-facing arc. The resulting finite graph is exactly
$A=\coprod_i A_i$. A loop contributes degree two at its vertex. Under the
geometric contract, each vertex has degree one or two; every connected
component is a path or a cycle. A path contributes one interval to $a_i$,
whereas a cycle contributes one attached boundary circle to $c_i$.

Supply, for each port, the oriented edge path followed by its central line
in a core graph. Endpoint images must agree after quotienting. Traversing
a cycle concatenates these paths, reversing and inverting a path when the
port is traversed backward. The result is a closed core word. For an
orientable base, the traversal must follow the boundary direction certified
by a single orientation of that base.

\begin{theorem}[Bounded marked-port extraction]
\label{thm:ports}
Given the correct quotient-vertex endpoint masks, coherent base-boundary
orientations, and core paths of total word length $P$, the numbers $a_i,c_i$
and all circle-image words can be constructed without visiting any unmarked
base cell. The finite graph has at most $2p$ vertices and $p$ edges. Each port
and each letter of its core path is traversed a bounded number of times.
The resulting interface and its verification have polynomial bit cost in
$p+P$ and the bit length of the supplied labels and profiles.
\end{theorem}
\begin{proof}
Decode the masks to a partition of $2p$ occurrences. Build the graph and
check ownership of vertices, degree one or two, and matching core images.
A traversal of each component distinguishes its two degree-one endpoints
from the degree-two cycle case. Record one interval anchor for a path.
For a cycle, concatenate its core paths in traversal order. Agreement at
quotient endpoints proves contiguity and closure. Each edge occurs once in
its component traversal, so the total emitted word length is at most $P$.
The bit masks have the bounds in Lemma~\ref{lem:masks}; sorting, graph
traversal, and exact integer comparison are polynomial in the stated input
size. No base genus or normal multiplicity is a traversal bound.
\end{proof}

\paragraph{The eight-port disc example.}
The maintained parallelity article describes a disc base supplied by an
index-two product handle, with eight core-facing arcs and no selected
index-zero or index-one handles \cite{repo-parallelity}. Its attachment
subset is a circle, so $a=0$, $c=b=1$, and $\chi(A)=0$. Treating it as eight
intervals would both alter Euler characteristic and lose its circle relation.
The delivered eight-port test exercises this combinatorics, but is not a
replay of that native triangulated example.

\paragraph{Exact scope of the implemented adapter.}
\code{ports.py} implements the mask decoder and the finite port graph
construction. It checks local directions along each attachment component.
It does not invoke the native weighted-orbit engine, recover source core
paths, or certify relative orientations of different boundary circles.
Those are the remaining inputs to a source adapter, not consequences of a
port histogram or a successful algebraic replay.

\section{Twisted attachments: torsion is tightly controlled}
\label{sec:twisting}

First perform the saturated orientable quotients and write
\begin{equation}
 G=H_1(K;\Z)/\langle s_i:i\in\satO\rangle.
 \label{eq:G}
\end{equation}
For the remaining saturated nonorientable bases, put
\[
 Q=(G\oplus\Z\langle h_1,\ldots,h_t\rangle)
       /\langle s_j+2h_j:1\leq j\leq t\rangle.
\]
Thus $Q=\coker C$ in the notation of Section~\ref{sec:condensation}.

\begin{theorem}[Twisted-attachment exact sequence]
\label{thm:twisting}
There is a natural short exact sequence
\begin{equation}
 0\longrightarrow G\longrightarrow Q\longrightarrow(\Z/2)^t
 \longrightarrow0.
 \label{eq:extension}
\end{equation}
Consequently:
\begin{enumerate}
\item The odd-primary torsion subgroups of $G$ and $Q$ agree under inclusion.
\item If $e$ is the exponent of $\Tor G$, with $e=1$ when $G$ is torsion-free,
then the exponent of $\Tor Q$ divides $2e$.
\item If $G\cong\Z^r$, and $\overline S$ has columns $s_j\bmod2$ in
$G/2G$, then
\begin{equation}
 Q\cong\Z^r\oplus(\Z/2)^{\,t-\rank_{\F}\overline S}.
 \label{eq:twist-rank}
\end{equation}
\end{enumerate}
In particular, $t$ twisted attachments do not produce a factor $2^t$ in the
\emph{torsion exponent} bound.
\end{theorem}

\begin{proof}
An element of $G$ that vanishes in $Q$ is an integer combination of the
relations $s_j+2h_j$. Comparing the independent $h_j$ coordinates in
$G\oplus\Z^t$ shows that every coefficient in this combination is zero.
Hence $G\to Q$ is injective. Modulo $G$, the generators $h_j$ have the
independent relations $2h_j=0$, proving \eqref{eq:extension}.

An element of odd prime-power order in $Q$ maps to zero in the elementary
$2$-group quotient, hence lies in $G$. Conversely, odd-primary torsion of
$G$ remains torsion of the same order by injectivity. This proves the first
claim. For any torsion element $q\in Q$, the element $2q$ lies in $G$ and is
torsion there. Thus $2eq=0$, proving the second claim.

When $G$ is free, tensoring \eqref{eq:extension} with $\Q$ gives free rank
$r$ for $Q$, and the second claim says its torsion is elementary $2$-torsion.
The presentation has $r+t$ generators and $t$ relations. Modulo two its
relation matrix has rank $\rank_{\F}\overline S$, because the lower block
$2I_t$ vanishes. Therefore
\[
 \dim_{\F}(Q/2Q)=r+t-\rank_{\F}\overline S.
\]
Subtracting the contribution $r$ from its free part gives
\eqref{eq:twist-rank}.
\end{proof}

\begin{example}[Large order is not large exponent]
Take $G=\Z/3\langle z\rangle$ and $s_j=z$ for all $j$. Then
\[
 Q\cong\Z/6\oplus(\Z/2)^{t-1}\qquad(t\geq1).
\]
Its torsion order is $3\cdot2^t$, but its exponent is always $6$.
The test suite checks the five-attachment instance and 700 random free-core
instances of \eqref{eq:twist-rank}.
\end{example}

\subsection{Canonical torsion without expanding repeated factors}
After a Smith computation on $C$, let its positive nonunit diagonal entries
be $d_1\mid\cdots\mid d_r$. The full torsion also contains $\mu$ copies of
$\Z/2$, where $\mu$ may be a very large binary integer. Returning a list of
$\mu$ entries would undo the compression.

\begin{proposition}[Run-length invariant factors]
\label{prop:torsion-runs}
The invariant factors of
$\bigoplus_j\Z/d_j\oplus(\Z/2)^\mu$ can be represented by $O(r+1)$
constant-value runs and computed in polynomial bit time in the input
diagonal and $\bits(\mu)$. No integer factorization is required.
\end{proposition}
\begin{proof}
Write $d_j=2^{a_j}u_j$ with $u_j$ odd. The $a_j$ are nondecreasing and the
$u_j$ form a divisibility chain. Remove trivial odd factors. Merge $\mu$
additional exponent-one terms into the ordered positive $a_j$, retaining
multiplicities in binary. Align this exponent sequence and the odd-factor
sequence on the right, padding the shorter on the left with zero exponents
or odd factor one. Multiply aligned factors.

For each prime, this is precisely the nondecreasing primary invariant-factor
sequence padded on the left by trivial factors. Their termwise products
are therefore the integral invariant factors. Only the $2$-adic valuations
are explicitly extracted; the odd factors remain whole integers.
The exponent sequence has at most $r+1$ runs and the odd sequence at most
$r+1$. Merging run endpoints uses a linear number of run operations and
binary arithmetic. No run is expanded according to its multiplicity.
\end{proof}

\section{Incidence bounds and normal-multiplicity safeguards}
\label{sec:torsion-bound}

A small number of handles alone does not bound the bit size of their
attaching maps. This distinction matters both for complexity and for
mathematical statements about torsion.

\begin{proposition}[Torsion from actual matrix incidence]
\label{prop:mass}
For an integer matrix $C$, set $S=\sum_{ij}|C_{ij}|$. Then
\[
 |\Tor(\coker C)|\leq 3^{S/3}.
\]
For the full compressed group in Theorem~\ref{thm:condensation},
\[
 |\Tor H_1(X)|\leq 2^\mu3^{S/3}.
\]
The latter is a mathematical bound, not a request to materialize $2^\mu$.
\end{proposition}
\begin{proof}
If the rank is zero, the first torsion group is trivial. Otherwise, let its
rank be $r$. Smith theory identifies the torsion order with the greatest
common divisor of the nonzero $r\times r$ minors. It is bounded above by
any nonzero such minor. Hadamard's inequality bounds that determinant by
the product of the Euclidean norms of its columns, and hence by their full
column $\ell^1$ norms $q_1,\ldots,q_r$. Each $q_j$ is a positive integer,
$\sum q_j\leq S$, and $q\leq3^{q/3}$ for every positive integer $q$.
The last inequality is immediate for $q=1,2,3$ and follows from the decrease
of $(\log x)/x$ for $x\geq3$. Multiplying proves the first bound. The isolated
$\Z/2$ factors supply the second.
\end{proof}

\begin{corollary}[Linear-incidence interfaces]
\label{cor:linear-incidence}
Suppose a family of decompositions has total core two-cell word length and
total attached-circle word length $O(T)$, and has $O(T)$ saturated twisted
bases. Then the torsion order of every component meeting the core is
$2^{O(T)}$, independently of all base genera and normal multiplicities.
An unattached interval-bundle component has torsion order at most two.
\end{corollary}
\begin{proof}
Forest collapse deletes edge coordinates. The total absolute mass of $R$
and the boundary-sum columns is no larger than the total input word length.
The lower block contributes $2t$. Thus $S=O(T)$. The matrix splits into
blocks according to the connected components of $\Gamma$; an unattached
bundle contributes no block and has the surface homology already computed.
Apply Proposition~\ref{prop:mass} separately to each component.
\end{proof}

\paragraph{What this says about the native source charts.}
The current repository supplies $O(T)$ core and port records, with bounded
total source-sector incidence \cite{repo-parallelity}. To use
Corollary~\ref{cor:linear-incidence} with that exact constant $T$, an adapter
must verify that its core graph and port paths have \emph{linear total word
length}, not merely that their number of records is linear. The present
package does not authenticate that refinement. A record containing a binary
power of a long winding cannot be charged as a single letter. With general
coefficient bit length $B_C$ and $q$ matrix rows, the safe alternative is the
usual determinant estimate $\log|\Tor\coker C|=O(q(B_C+\log(q+1)))$.

\subsection{Closed one-sided components cannot proliferate arbitrarily}
\begin{proposition}
\label{prop:one-sided}
If a compact orientable three-manifold $M$ contains $s$ pairwise disjoint
closed connected nonorientable embedded surfaces, then
$s\leq b_1(M;\F)$. Consequently, the counter $\mu$ in an actual cut of a
fixed source $M$ is at most $b_1(M;\F)$.
\end{proposition}
\begin{proof}
Each surface is one-sided. Choose an orientation-reversing loop on it.
Push the loop slightly into its normal interval bundle: it starts on one
local side and ends on the other. Closing it by a short fibre segment
produces a loop that meets this surface once modulo two. The loop lies in
that surface's bundle and is disjoint from all the other surfaces.

Intersection parity with each closed surface defines a class in
$H^1(M;\F)$. Evaluating these classes on the constructed loops yields the
identity matrix. The classes are linearly independent, proving the bound.
For the second assertion, take the central closed base surfaces of the
unattached nonorientable bundles. A cut is realized as a submanifold of
$M$, so these surfaces are also pairwise disjoint in $M$.
\end{proof}

For a triangulation with $T$ tetrahedra, the crude bound
$b_1(M;\F)\leq6T$ follows from the number of source edges. Together with a
verified linear-incidence adapter, this would bound the \emph{total}
torsion order over all cut components by $2^{O(T)}$. The very large
multiplicity tests in our abstract model therefore must not be portrayed as
normal cuts of a fixed small triangulation; they test the encoded algebraic
contract beyond that geometric restriction.

\section{A stronger shortcut for certified descendants in \texorpdfstring{$S^3$}{S3}}
\label{sec:s3}

The generic API accepts a compact orientable triangulation with one torus
boundary. This is a weaker promise than being the exterior of the input knot
in $S^3$. The distinction permits a useful optimization only after source
provenance is carried explicitly.

\begin{theorem}[Boundary-determined integral homology in $S^3$]
\label{thm:s3}
Let $X\subset S^3$ be a connected compact locally flat PL codimension-zero
submanifold with nonempty boundary. Let $b$ be the number of its boundary
components and let $G_\partial$ be their total genus. Then
\begin{equation}
 (H_0,H_1,H_2,H_3)(X;\Z)
 =\bigl(\Z,\Z^{G_\partial},\Z^{b-1},0\bigr).
 \label{eq:s3-homology}
\end{equation}
In particular, all these integral homology groups are torsion-free.
\end{theorem}

\begin{proof}
Let $\Sigma=\partial X$, a disjoint union of $b$ closed orientable surfaces.
Alexander duality gives
\[
 \widetilde H_0(S^3-\Sigma;\Z)\cong H^2(\Sigma;\Z)\cong\Z^b.
\]
Thus $S^3-\Sigma$ has $b+1$ connected regions. The connected interior of $X$
is exactly one region: it is open and closed in $S^3-\Sigma$ and has no
additional boundary there. Hence $S^3-X$ has $b$ components.

Let $Y$ be the closure of $S^3-X$. Its collar identifies $Y$ up to homotopy
with $S^3-X$, and identifies $S^3-Y=\operatorname{int}X$ up to homotopy with
$X$. Apply Alexander duality to the compact locally contractible space $Y$.
This gives
\[
 H_2(X;\Z)\cong\widetilde H^0(S^3-X;\Z)\cong\Z^{b-1},
 \qquad H_1(X;\Z)\cong H^1(S^3-X;\Z).
\]
The right-hand group in degree one is torsion-free by the universal
coefficient theorem. It is finitely generated, since a collar makes the
complement homotopy equivalent to a compact manifold. Finally, $H_3(X)=0$
and doubling $X$ along its boundary gives
$2\chi(X)=\chi(\partial X)=2b-2G_\partial$. If $r=b_1(X)$, the already
computed groups give $\chi(X)=1-r+(b-1)=b-r$. Therefore $r=G_\partial$.
These duality statements are standard; see \cite{hatcher} and the related
embedding obstructions in \cite{tonkonog}.
\end{proof}

\begin{corollary}[A source-aware sphere terminal]
\label{cor:ball}
Under the hypotheses of Theorem~\ref{thm:s3}, $b=1$ and $G_\partial=0$
imply that $X$ is a three-ball. More generally, if $G_\partial=0$, then $X$
is homeomorphic to $S^3$ with the interiors of $b$ disjoint three-balls removed.
\end{corollary}
\begin{proof}
Each boundary component is a locally flat two-sphere. The Schoenflies
theorem says that the closures of both complementary domains of such a
sphere are three-balls \cite[Theorem 5]{brown}. There are $b$ complementary
components of $X$, and $b$ boundary components. Every complementary
component meets at least one boundary component; thus each meets exactly
one sphere and is a ball. Removing their interiors gives the assertion.
For $b=1$, either complementary closure is a ball.
\end{proof}

\paragraph{Why homology alone is not the terminal.}
The conclusion is a theorem about an embedded manifold, not an inference
from a Betti vector. A punctured lens space has one sphere boundary and can
have nontrivial torsion; a punctured integral homology sphere can have exactly
the homology of a ball without being a ball. Neither is a counterexample to
Corollary~\ref{cor:ball}, because neither can satisfy its $S^3$ embedding
premise in the nontrivial case. The existing generic sphere-boundary nonball
control should therefore be retained, not removed.

\subsection{Primitive relations and a parity obstruction}
\begin{corollary}[Necessary interface conditions in $S^3$]
\label{cor:primitive}
Assume the full decomposition \eqref{eq:decomp}, including a regularized
core $K$, lies in $S^3$. Then:
\begin{enumerate}
\item $\mu=0$ and $H_1(K;\Z)$ is free;
\item the subgroup generated by the saturated orientable sums $s_i$ is a
primitive sublattice of $H_1(K;\Z)$;
\item writing $G$ as in \eqref{eq:G}, the saturated nonorientable sums are
linearly independent in $G/2G$. In particular $t\leq\rank G$.
\end{enumerate}
If $X$ is connected, then
\begin{equation}
 G_\partial=\rank G+\beta.
 \label{eq:s3-rank}
\end{equation}
\end{corollary}
\begin{proof}
Theorem~\ref{thm:s3} applies to every connected component of $X$ and $K$.
Thus their first homologies are free. Formula~\eqref{eq:h1} forces $\mu=0$.
Theorem~\ref{thm:twisting} injects $G$ into $Q$, and $Q$ is a direct summand
of the torsion-free group $H_1(X)$. Therefore $G$ is torsion-free. A quotient
of a free abelian group by a subgroup is torsion-free exactly when that
subgroup is primitive, giving the second assertion.

Now $G$ is free. Formula~\eqref{eq:twist-rank} and torsion-freeness of $Q$
force $\rank_{\F}\overline S=t$. Rationally, $\rank Q=\rank G$; combine
this with \eqref{eq:h1} and \eqref{eq:s3-homology} for the last identity.
\end{proof}

These are necessary conditions, not an embedding recognizer. Passing a
primitivity or parity test does not establish that the interface came from
$S^3$. Conversely, torsion discovered under a claimed $S^3$ ancestry signals
a defective ancestry, geometric adapter, or algebraic computation. It is
not, by itself, a knottedness obstruction: nontrivial knot exteriors also
embed in $S^3$.

\subsection{How provenance can be preserved}
An implementation can retain an embedding-ancestry DAG rather than expand
coordinates in $S^3$ at every stage. Its root must bind a certified
construction of the input knot exterior to the original diagram. Cutting
out an open regular neighborhood of a properly embedded surface preserves
the embedding by inclusion, and selecting components preserves it as well.
A verified homeomorphic change of representation also preserves the
existence of an embedding, by composition with the inverse homeomorphism.

An arbitrary quotient, filling, or gluing is not covered by these rules.
Neither a hash nor a Boolean field named \code{certified} proves the root
premise or a transition. This delivery includes no checker for such a DAG.
The deliberately named \code{s3\_boundary\_consequences} computes only the
consequent of Theorem~\ref{thm:s3}; its returned record explicitly says that
the embedding hypothesis has not been verified by that function.

Once the ancestry and the boundary profiles are independently replayed,
\eqref{eq:s3-homology} costs only binary arithmetic per profile. It can
replace an interior matrix computation for Betti numbers or the predicate
$b_1=0$. It does not supply the actual nonseparating surface, a boundary
pattern, or the essentiality test needed by the hierarchy.

\section{Compact cohomology witnesses}
\label{sec:cohomology}

A hierarchy needs more than a nonzero Betti number. We can nevertheless
separate finding an algebraic functional from lifting it to a geometric
cutting surface.

\begin{theorem}[A small functional in the condensed coordinates]
\label{thm:functional}
Given the interface data, if $H^1(X;\Z)\ne0$, then a nonzero integral
functional on the presentation in Theorem~\ref{thm:condensation} has a
description polynomial in the encoded input size. It can be found with a
polynomial-time integer normal-form algorithm. The description uses at most
one binary-indexed bulk generator or an ordinary vector on the $d+t$
interacting generators. It need not enumerate a surface homology basis.
\end{theorem}
\begin{proof}
If $\beta>0$, choose a nonzero coordinate functional on one of the free
summands recorded in \eqref{eq:bulk}. A directory consisting of the
$r_\Gamma$ graph generators followed by the surface-profile generators
locates the chosen generator using a profile index, copy index, and local
basis index, all in binary. If desired, choose the first nonempty family,
which does not even require an index as large as the whole rank.

Otherwise, $H^1(X;\Z)$ is the integer left kernel of $C$; finite torsion
contributes no homomorphism to $\Z$. A polynomial-time integer normal-form
algorithm gives a nonzero vector in this kernel whenever its rational
nullity is positive \cite{kannan-bachem}. Its coefficient bit lengths are
polynomially bounded. Dividing the vector by the greatest common divisor
of its entries makes it primitive without changing its support.

For completeness, the symbolic local coordinates can be interpreted in the
standard surface presentations without expanding them. In an unsaturated
base, eliminate one unattached boundary generator $d_b$ and assign its
functional value using the single surface relation. If one crosscap
coordinate is set to one, only $d_b$ needs the compensating value $-2$;
if one free boundary generator is selected, only $d_b$ needs $-1$.
In a saturated nonorientable base use
$x_1=h-\sum_{j\geq2}y_j$, $x_j=y_j$. A single $y_j$ coordinate therefore
requires only two nonzero original crosscap values. Orientable handle
coordinates are already independent. A graph functional is represented on
one nonforest edge of $\Gamma$. These are constant-support recipes with
binary indices, or ordinary vectors of length $d+t$.
\end{proof}

\paragraph{The remaining geometric lift is real work.}
Theorem~\ref{thm:functional} does not identify an explicit source simplicial
cocycle, let alone a normal surface. A native lift must transport the
canonical base coordinates through the actual affine polygon identifications
and respect the core inclusion. This requires certified chain maps or an
equivalent source-level evaluation mechanism. In particular, reporting
``$b_1>0$'' and invoking a surface-construction theorem without supplying its
required cochain is not an implementation.

A useful next target is a \emph{port-preserving cocycle-lift theorem}: from a
replayed affine base/core packet and one of the above functionals, construct
a source-evaluable integral cocycle in polynomial encoded time. Its input
and output must include the actual surface chart maps. This article isolates
that question but does not claim its solution.

\section{Bit complexity and its role in the quasi-polynomial goal}
\label{sec:complexity}

\subsection{The local bound}
Let $N$ be the complete bit length of the supplied core, profile, and
interface data, including all explicit attaching words. Let $h$ be the
number of core cells and $p$ the number of raw ports. The interface has at
most $p$ circle components and at most $p$ interval components. In
particular $o+t\leq p$ and $d\leq h$.

\begin{theorem}[Polynomial encoded homology]
\label{thm:complexity}
From a supplied verified interface, constructing \eqref{eq:matrix},
\eqref{eq:bulk}, \eqref{eq:h0}, and \eqref{eq:euler} has polynomial bit cost
in $N$. The matrix has $O(h+p)$ rows and columns. A polynomial-time integer
Smith algorithm, followed by Proposition~\ref{prop:torsion-runs}, computes
a complete run-length description of $H_1(X;\Z)$ in polynomial bit time.
No dependence on the numerical genus or multiplicity is hidden in an
output-length requirement.
\end{theorem}
\begin{proof}
A spanning forest, core word checks, and attachment graph construction
have polynomial cost. Every matrix coefficient is an exponent sum over
supplied words, so its bit length is bounded by the word encoding size.
There are at most $h+p$ generators and $h+p$ relations up to fixed constants.
Dense assembly is therefore polynomial even without sparse optimization.
The $q$ profile contributions require $O(q)$ arithmetic operations on binary
integers of polynomially bounded bit length. Products such as $m_i\lambda_i$
have bit length at most the sum of the factor bit lengths plus one.

Kannan--Bachem provides a polynomial integer normal-form algorithm, including
controlled transformation matrices \cite{kannan-bachem}. The small Smith
output has polynomial size. The run-length merge in
Proposition~\ref{prop:torsion-runs} retains binary multiplicities. Free
summands are likewise represented by their rank. Hence the output is a
compressed isomorphism type rather than an exponentially long basis.
\end{proof}

The theorem begins \emph{after} source interface extraction. The actual
repository's weighted discovery and replay costs must be added with their
stated trace/partition accounting \cite{repo-parallelity}. A polynomial
postprocessor cannot retroactively prove a bound on a preceding discovery
algorithm whose largest intermediate encoding has not been bounded.

\subsection{The reference Smith kernel is not a Kannan--Bachem implementation}
The delivered optional Smith kernel uses elementary Euclidean row and column
operations, emitting a trace. It is exact and terminating. To see
termination, at a pivot clear its column and row by division with remainder.
Every nonzero remainder that becomes the new pivot has smaller absolute
value. Once the pivot row and column are clear, an entry in the remaining
block not divisible by the pivot is moved into the pivot row; another
remainder again decreases the pivot. Eventually the pivot divides the
remaining block and the algorithm advances to a smaller submatrix.
Previously settled pivots remain isolated, and later integral operations
preserve divisibility by them.

This argument is not a bound on all intermediate coefficient sizes. We do
not claim that this particular implementation has Kannan--Bachem's worst-case
polynomial bit complexity. The proved polynomial \emph{algorithmic existence}
result uses that published normal-form backend; the delivered prototype's
normal-form cost is measured and replayed. The polynomial assembly stage
can also be used independently of any Smith computation.

\subsection{What would close the global accounting}
\begin{proposition}[A sufficient accounting criterion]
\label{prop:qp}
Suppose a complete, correct hierarchy recognizer on an $n$-crossing diagram
has at most $L(g+1)^L$ iterations, each iteration costs at most $P(N)$ bit
operations, and all preprocessing is bounded by the same total scale.
If $P$ is a fixed polynomial, $L=O(\log n)$, $g=n^{O(1)}$, and
$\log N=O((\log n)^2)$, its total running time is
$2^{O((\log n)^2)}$. More generally a sufficient bound for
$2^{O((\log n)^3)}$ is
\[
 \log P(N)+\log(L+1)+L\log(g+1)=O((\log n)^3).
\]
\end{proposition}
\begin{proof}
Multiply the iteration bound by the largest per-iteration cost and take
logarithms. Under the first hypotheses, both $\log P(N)$ and
$L\log(g+1)$ are $O((\log n)^2)$; the extra logarithmic factor in $L$ is
smaller. The second assertion is the same calculation without specialization.
\end{proof}

This is sufficient accounting, not a proof that the actual hierarchy has
these bounds. The present contribution removes numerical base genus and
multiplicity from one matrix dimension. It does not establish bounded
patterned-core encodings throughout every correction, a logarithmic depth,
or the completeness and cost of the constructive hierarchy operations.
The current complete recognizer's fallback is not replaced.

\section{Implementation and verification artifacts}
\label{sec:implementation}

\subsection{The module boundary}
The package is an additive research component, intended for placement under
\path{Topology/UnknotRecognition/research/interface_homology/}. No native
recognition dispatch is modified. Its main modules are:
\begin{description}[style=nextline]
\item[\code{model.py} and \code{condense.py}]
Strict integer validation, closed edge words, binary profiles, forest
collapse, the attachment graph, the matrix $C$, and all bulk counters.
\item[\code{ports.py} and \code{replay.py}]
Marked endpoint decoding and finite port assembly; independent replay of the
algebraic condensation using graph traversal rather than the producer's
disjoint-set implementation.
\item[\code{smith.py} and \code{boundary.py}]
The small exact Smith kernel, elementary-operation replay, binary torsion
run merging, a mod-two rank utility, and explicitly conditional boundary
consequences for an $S^3$ descendant.
\end{description}
SymPy is a validation dependency only. The main package and benchmark kernel
use the Python standard library. No third-party mathematical paper or font
file is redistributed.

\subsection{Certificate meanings}
A condensation packet records the source-model SHA-256 fingerprint, chosen
core forest, surviving edges, component labels, the matrix, saturated-base
indices, graph counts, and bulk counters. Replay checks a spanning forest
by independent component traversal, extracts exponent sums independently,
and recomputes all formulas. The hash binds an algebraic model. It is not
an assertion that the model came from the diagram.

A Smith trace allows only row/column swaps, adding an integer multiple of
one row or column to a \emph{different} row or column, and row negation.
Every operation is unimodular. Replay applies the supplied operations and
checks that the result is diagonal, has positive nonzero entries in a
prefix, and satisfies successive divisibility. These checks prove equality
of the presented abelian groups without relying on the pivot-selection code.

Operation and bit caps produce \code{WorkLimit}; they do not produce a
partial homology group or a topological verdict. The replay operation cap
limits the supplied trace. For untrusted deployment, external input-size
and elapsed-time limits are still appropriate; the research API is not
advertised as a hardened network service.

\subsection{Binary serialization}
Python environments can limit conversion of enormous integers to decimal
strings. The delivered wire codec encodes integers as signed canonical
hexadecimal strings under an explicit tag. It rejects ambiguous forms such
as negative zero, leading zeroes, or a prefixed hexadecimal spelling inside
that tag. Booleans remain booleans and are rejected where an integer is
required. The fingerprint uses this codec, so values such as $2^{16384}$ do
not require changing a process-wide decimal-conversion limit.

\subsection{What the tests do and do not certify}
The independent expanded oracle writes down the standard abelian surface
presentations \emph{before} the eliminations in Theorem~\ref{thm:condensation}.
It uses its own breadth-first core forest. Small expanded matrices are
compared using SymPy's Smith normal form \cite{sympy}. A second family
checks the reference Smith kernel against all determinantal divisors of
small matrices. These are independent finite algebraic checks, not proofs
of a source triangulation's geometry.

The current delivery has no Regina installation, does not execute the
maintained repository suite, and does not replay the native cut-cell or
parallelity certificates. It supplies neither a Lean formalization nor a
full source-to-interface producer. These limitations are recorded in the
machine-readable validation output as well as in the integration notes.

\section{Completed experiments and reproducibility}
\label{sec:experiments}

\subsection{Validation families}
The delivered run passes all 25 unittest methods, with no failures, errors,
or skips. Within those methods, six seeded families contain 4,800 cases:
\begin{center}
\begin{tabular}{@{}lr@{}}
\toprule
Independent comparison or property & Cases\\
\midrule
Expanded surface presentation versus condensation & 1,200\\
Reference Smith diagonal versus SymPy & 1,800\\
All determinantal divisors of small matrices & 200\\
Binary repeated-$2$ invariant-factor merging & 500\\
Free-core twisted-attachment formula & 700\\
Oriented marked-port cycles & 400\\
\bottomrule
\end{tabular}
\end{center}
The six random seeds are the integers from \code{2026101201} through
\code{2026101206}. Deterministic tests additionally cover empty matrices and
cores, disconnected cores, interval cycles, both signs of boundary maps,
M\"obius attachments, closed surfaces, malformed models, forged packets,
resource caps, and 16,385-bit parameters. The summary in
\code{results/validation.json} records the environment and source hashes.
These counts describe this delivery, not a rerun of any historical native
checkpoint.

\subsection{Paired complete homology-kernel measurements}
Both timing arms start with the same abstract interface model and finish
with a serialized $H_0/H_1$ answer. The condensed arm assembles and replays
its packet, computes Smith form of the small matrix, replays that trace, and
merges the binary factors. The expanded arm builds the literal standard
surface presentation, runs the \emph{same} Smith kernel, replays its trace,
and serializes the same invariant-factor answer. Thus the comparison
isolates presentation compression rather than a change of normal-form
library. The condensed arm includes its additional condensation replay.

There are eleven cases, one warmup round, and nine measured paired rounds.
A third arm repeats the condensed computation as an A/A control. The three
arms are shuffled each round with seed \code{2026101211}; garbage collection
is performed outside each timer. All 330 warmup and measured samples finish,
and all answers agree. Of these, 297 are measured samples. Tables report the
median of the \emph{paired} expanded/condensed ratios, not the ratio of two
independently rounded medians. Values greater than one favor condensation.

\begin{table}[ht]
\centering\small
\begin{tabular}{@{}lrrrr@{}}
\toprule
Case & Condensed (ms) & Expanded (ms) & Paired ratio & A/A\\
\midrule
$g=0$ & 0.284 & 0.182 & 0.647 & 1.065 \\
$g=1$ & 0.290 & 0.199 & 0.677 & 1.038 \\
$g=8$ & 0.302 & 0.260 & 0.895 & 0.954 \\
$g=64$ & 0.299 & 0.777 & 2.507 & 0.995 \\
$g=512$ & 0.283 & 4.564 & 16.123 & 0.983 \\
$g=4096$ & 0.284 & 33.603 & 116.599 & 0.974 \\
$m=1$ & 0.235 & 0.141 & 0.625 & 1.044 \\
$m=8$ & 0.235 & 0.680 & 2.887 & 0.980 \\
$m=32$ & 0.239 & 24.992 & 108.174 & 1.013 \\
core only tiny & 0.215 & 0.113 & 0.517 & 0.980 \\
core only relations & 1.805 & 1.452 & 0.796 & 1.038 \\
\bottomrule
\end{tabular}
\caption{Complete abstract $H_1$ kernels, including exact Smith replay and output. Nine measured rounds; medians and paired ratios. No source cut or knot-recognition time is included.}
\label{tab:timings}
\end{table}


The saturated-genus family attaches an orientable genus-$g$, two-boundary
base to a one-loop core with boundary images $z,-z$. Its compressed matrix
is always $1\times1$, while the expanded presentation has $2g+4$ rows and
three columns; one row accounts for the attachment-graph cycle. The exact
answer is $H_1=\Z^{2g+2}$.

The unattached nonorientable family has $m$ copies of a closed three-crosscap
base. Its exact first homology is $\Z^{2m}\oplus(\Z/2)^m$. It is a legitimate
abstract bundle model but, by Proposition~\ref{prop:one-sided}, large $m$
is not a cut of a fixed small source. The core-only controls have no bulk
surface data and expose the extra packet/replay overhead.

The observed speedups are local. In particular, the genus-4096 case is a
homology-model measurement, not a 4096-crossing knot or a measured native
normal cut. Tiny cases regress, and even an excellent matrix-stage ratio
says nothing about a source extractor that dominates the total runtime.

\subsection{Large binary parameters without an expanded baseline}
For $N=2^e$, we also process $N$ unattached copies of a closed nonorientable
base with $N+1$ crosscaps. The exact answer is
\[
 H_0=\Z^N,\qquad H_1=\Z^{N^2}\oplus(\Z/2)^N.
\]
The interacting matrix is empty. Source-model hashing, packet replay, the
homology result, and its serialization are included. The first of ten runs
is treated as warmup.

\begin{table}[ht]
\centering
\begin{tabular}{@{}rrrr@{}}
\toprule
$e$ in $N=2^e$ & Matrix & Serialized answer bytes & Median ms\\
\midrule
64 & $0\times0$ & 157 & 0.058 \\
256 & $0\times0$ & 349 & 0.055 \\
1,024 & $0\times0$ & 1,117 & 0.062 \\
4,096 & $0\times0$ & 4,189 & 0.107 \\
16,384 & $0\times0$ & 16,477 & 0.272 \\
\bottomrule
\end{tabular}
\caption{Binary stress tests, with genus/crosscaps and multiplicity both large. No expanded baseline or timing ratio is claimed.}
\label{tab:binary}
\end{table}


The expanded arm is deliberately not run at these sizes, and no timing
ratio against it is asserted. The purpose is to test that construction,
replay, and output serialization really use bit length rather than numerical
multiplicity. This is not evidence that a fixed-size normal triangulation
realizes those particular large profiles.

\subsection{Reproduction commands}
From the package root:
\begin{lstlisting}
python validate.py
python benchmark.py
python make_tables.py
sh build.sh
\end{lstlisting}
The independent tests require SymPy; the implementation and benchmark
kernel do not. The completed run used Python and SymPy versions recorded
in the JSON files. Timing results are specific to this environment and
should be regenerated under the project's normal native benchmark controls
before selecting a production threshold.

\section{Integration plan and corrections to the research interface}
\label{sec:integration}

\subsection{A narrow, testable native adapter}
The proposed adapter should consume a successfully replayed
\code{normal-cut-parallelity-v1} packet and produce four things: a core
one-skeleton with two-cell words; the endpoint equivalence classes of all
core ports; a central core path for each port; and relative boundary
orientation data for orientable bases. Its source fingerprint must match
the original triangulation and unscaled normal vector.

The existing base profiles already supply the genus or crosscap count,
boundary count, multiplicity, and port ownership. They do not by themselves
supply the ordered boundary cycles or their signed maps. The finite adapter
in \code{ports.py} starts precisely after those source-level inputs have
been obtained. It should not be connected to a guessed geometric map or a
profile-only approximation.

Validation should proceed on small native cuts by comparing the resulting
integral group with an independent expanded-cut homology calculation,
including torsion, not just ranks over two finite fields. Relabelled source
triangulations, one-sided multiples, high-degree edge links, and the
index-two-only disc example should be included. Tests must also disable the
producer while replaying saved source evidence, so that an algebraic
agreement cannot mask identical geometry mistakes on both sides.

\subsection{A separate provenance-aware boundary path}
The $S^3$ optimization has a different adapter: a replayed diagram-to-exterior
root, a cut/homeomorphism ancestry, and componentwise boundary profiles.
Its conclusion is exact integral homology by Theorem~\ref{thm:s3}; a
single-sphere component additionally has a ball conclusion by
Corollary~\ref{cor:ball}. The generic triangulation path remains available
and retains the sphere-boundary nonball counterexample.

The algebraic path remains valuable even when the boundary shortcut applies:
it can check necessary interface identities, retain core inclusion data,
and prepare cohomology functionals. It is not necessary to compute Smith
form merely to rediscover a torsion-free Betti vector already implied by
verified $S^3$ provenance.

\subsection{Documentary and specification corrections}
The audited root README still uses the broad phrase ``bounded-core geometry''
in its earlier boundary-constructor checkpoint. The newer
\code{parallelity.tex} proves a bounded \emph{inherited} core. A precise
status update should distinguish that achievement from the still-missing
admissible patterned core; it should not continue to list all bounded-core
construction as absent. Conversely, the statement that a bounded
\emph{admissible} core remains open is not contradicted by the inherited
core result.

The generic raw-cut article correctly keeps interior homology separate from
boundary data. That separation can be sharpened, not discarded: for inputs
with certified $S^3$ ancestry, boundary data already determines all integral
homology. For unrestricted torus-boundary triangulations, it does not.
No existing nonball test should be reclassified as incorrect on the basis
of the stronger, currently unverified source promise.

Finally, a source wrapper should avoid the phrase ``polynomial homology
implementation'' without specifying its Smith backend and the cost of its
geometric extractor. The mathematical polynomial-time theorem, the tested
Euclidean prototype, and the native weighted trace bounds are three
different assertions.

\section{Further research questions}
\label{sec:future}

\subsection{Source-bound port cycles and relative orientations}
Can the existing affine source charts emit all core words and oriented
attachment circles with a proved linear total incidence bound in $T$?
The endpoint-mask theorem reduces the component-classification part to one
marked vertex query and $O(p)$ finite graph data. The remaining issue is
source-certified paths and the relative signs on all boundary circles of
each orientable base. A useful first deliverable is an independently replayed
adapter with mutation tests for whole-circle orientation reversals.

\subsection{A verified \texorpdfstring{$S^3$}{S3} ancestry discipline}
Which current hierarchy operations preserve an explicit embedded descendant,
and which preserve only an abstract homeomorphism type? A small proof kernel
for root construction, cutting, component selection, and verified
homeomorphisms would make the boundary-only homology and sphere terminal
available without solving a new embedding problem at every stage. Operations
that fill or identify boundary pieces require separately proved rules.

\subsection{Source-evaluable cohomology and connected cutting surfaces}
Can the symbolic functionals of Theorem~\ref{thm:functional} be pulled back
through the affine bundle charts in polynomial bit time, with no expansion
of $g_i$ or $m_i$? The output should support evaluation on source edges and
independent checking of the cocycle equation. A subsequent step must produce
the primitive nonseparating surface required by the hierarchy, with its own
coordinate and construction bounds. Neither step follows from a rank count.

\subsection{Marked inclusion data beyond \texorpdfstring{$H_1$}{H1} as an abstract group}
Can one retain, at bounded-interface cost, the maps from the individual
boundary components into $H_1(X)$, intersection pairings, and the marked
peripheral classes? These data are stronger than an isomorphism type and
could guide a pattern-preserving simplification. The matrix $C$ already
retains the core inclusion, but its bulk direct-sum decomposition is not a
canonical boundary marking.

\subsection{Admissible patterned cores and repeated corrections}
The inherited $30T$-handle core is an important local bound. How can it be
converted to a uniform-type admissible structure while transporting every
boundary pattern, and how does that encoding behave after a weak reduction
or a discarded hierarchy suffix? A satisfactory theorem must bound the
largest intermediate bit length, not just the number of handles after one
normal cut.

\subsection{Logarithmic depth and constructive balancing}
A balance theorem must control both the topology and its construction cost.
Can a Cheeger-region or multisurface argument bound hierarchy depth by
$O(\log n)$ while keeping all correction interfaces within the same encoded
budget? Proposition~\ref{prop:qp} identifies the accounting target, but no
balanced-region search or depth proof is established here.

\subsection{Fast exact algebra with portable evidence}
The reference Smith trace is easy to replay but can exhibit coefficient
growth. A production backend could emit bounded-size unimodular certificates
from a polynomial normal-form algorithm, or use modular methods followed by
exact reconstruction and replay. When the source is in $S^3$, the primitivity
and parity constraints of Corollary~\ref{cor:primitive} offer additional
checks, but do not remove the need to authenticate the geometric source.

\subsection{Native measurements rather than model extrapolation}
The next benchmark should time the complete chain from a source
triangulation and normal vector through geometric replay, interface
construction, homology, and output. It should include the existing expanded
cut oracle, repeated A/A controls, and cases where the direct boundary path
is enabled only by certified ancestry. The small-input regressions in the
present model measurements are a reason to measure dispatch policies, not
to make condensation the unconditional default.

\section{Conclusion}
The interacting integral homology of a core with surface bundles is much
smaller than the expanded surface homology suggests. Unsaturated bases
contribute only free ranks; saturated orientable bases add one signed
boundary-sum relation; saturated nonorientable bases add one relation
$2h+s=0$ and one generator. This yields a complete compressed homology
presentation with independent algebraic replay and precise bit-cost
separation from the source extractor.

For certified $S^3$ descendants, the stronger boundary-only route avoids even
that matrix when only homology ranks or a sphere terminal are needed. Its
premise must remain explicit. The primitive-lattice and parity consequences
provide additional consistency checks, not a substitute embedding theorem.
These local results advance the compressed hierarchy interface while leaving
the construction, admissibility, depth, and global search bounds as explicit
mathematical obligations.

\begin{thebibliography}{99}
\raggedright
\bibitem{repo-cells}
ProveIt contributors, \emph{Full three-dimensional normal cuts in affine cell
families}, \code{Topology/UnknotRecognition/synthesis/cut\_cells.tex}, audited
revision listed in Appendix~\ref{app:audit}, 2026.
\url{https://github.com/VladimirReshetnikov/ProveIt}.

\bibitem{repo-parallelity}
ProveIt contributors, \emph{Source parallelity in all handle dimensions and a
bounded inherited core}, \code{synthesis/parallelity.tex}, and the associated
\code{fastunknot/normal\_parallelity.py}, audited revision, 2026.
\url{https://github.com/VladimirReshetnikov/ProveIt}.

\bibitem{lackenby-cert}
Marc Lackenby, \emph{The efficient certification of knottedness and Thurston
norm}, arXiv:1604.00290v3, especially Section 9.4, Theorem 9.3.
\url{https://arxiv.org/html/1604.00290v3}.

\bibitem{lackenby-2026}
Marc Lackenby, \emph{Incompressible surfaces, hierarchies and unknot
recognition}, arXiv:2607.23350v1, 25 July 2026, especially Sections 6, 9,
and 10. \url{https://arxiv.org/abs/2607.23350}.

\bibitem{hatcher}
Allen Hatcher, \emph{Algebraic Topology}, Cambridge University Press, 2002.
Surface homology, Mayer--Vietoris, universal coefficients, and duality.
\url{https://pi.math.cornell.edu/~hatcher/AT/ATpage.html}.

\bibitem{brown}
Morton Brown, \emph{A proof of the generalized Schoenflies theorem},
Bulletin of the American Mathematical Society \textbf{66} (1960), 74--76,
Theorem 5.
\url{https://projecteuclid.org/journals/bulletin-of-the-american-mathematical-society/volume-66/issue-2/A-proof-of-the-generalized-Schoenflies-theorem/bams/1183523455.full}.

\bibitem{kannan-bachem}
Ravindran Kannan and Achim Bachem, \emph{Polynomial algorithms for computing
the Smith and Hermite normal forms of an integer matrix}, SIAM Journal on
Computing \textbf{8} (1979), 499--507.
\url{https://doi.org/10.1137/0208040}.

\bibitem{tonkonog}
Dmitry Tonkonog, \emph{Embedding 3-manifolds with boundary into closed
3-manifolds}, Topology and its Applications \textbf{158} (2011), 1157--1162.
\url{https://arxiv.org/abs/1003.3029}.

\bibitem{sympy}
SymPy developers, \emph{Matrix normal forms}, version 1.14.0 documentation,
accessed 12 October 2026.
\url{https://docs.sympy.org/latest/modules/matrices/normalforms.html}.
\end{thebibliography}

\appendix
\section{Source audit and proof-status inventory}
\label{app:audit}

\subsection{Repository source pins}
The audited commit is
\begin{center}\small\code{ed166080d7db14dd7a91d07788996f02947dd92a}.\end{center}
The following paths are relative to \code{Topology/UnknotRecognition/}.
These Git blob identifiers bind the inspected texts; they are not SHA-256
hashes of the newly delivered files.
\begin{description}[style=nextline]
\item[\code{README.md}]
\code{4fb929b5689d0bc20327e241cb775789660c8013}.
\item[\code{synthesis/cut\_cells.tex}]
\code{dc12b1a58aedaed78f00b63109b0aa5214832ca1}.
\item[\code{synthesis/parallelity.tex}]
\code{9e6698d589568dafc35d7a39df7ad6f74ee46ca0}.
\item[\code{fast/fastunknot/normal\_parallelity.py}]
\code{22cf15ffa11ca85e6a6c350f5d217f1649a01ea2}.
\end{description}
The initial part of \code{normal\_cut\_boundary.py} was also inspected, but
its entire implementation was not audited. Code search for a Schoenflies
reference returned no hit in the inspected repository; this is not an
exhaustive novelty search. No claim is made to have read or rerun the many
historical reports and archives.

\subsection{Claim classes}
\begin{longtable}{@{}>{\raggedright\arraybackslash}p{.31\textwidth}>{\raggedright\arraybackslash}p{.61\textwidth}@{}}
\toprule
Claim & Status in this delivery\\
\midrule
\endhead
Integral condensation and graph rank & Proved for the stated interface
contract; standard-library implementation and independent algebraic replay.\\
Marked endpoint extraction & Proved and implemented downstream of the
correct marked quotient and source paths; no native orbit invocation.\\
Twisted exact sequence and parity formula & Proved; 700 random free-core
checks plus deterministic torsion tests.\\
Canonical binary torsion runs & Proved and implemented; 500 expanded
comparisons and a 16,385-bit multiplicity case.\\
Polynomial-time homology theorem & Proved using a polynomial integer
normal-form algorithm; the delivered Euclidean Smith kernel is not claimed
to have that bit bound.\\
Boundary-only homology and sphere terminal & Classical consequences proved
under certified $S^3$ embedding; the delivered helper does not verify the
premise and is not a recognizer.\\
Source geometry and diagram ancestry & Adapter contracts and proposed proof
rules only; not implemented or authenticated in this package.\\
Global quasi-polynomial recognition & Sufficient cost criterion and specific
remaining questions; no complete algorithm or global bound claimed.\\
\bottomrule
\end{longtable}

\section{Minimal executable example}
\label{app:example}
The following example is an algebraic saturated M\"obius attachment, not a
knot input:
\begin{lstlisting}[language=Python]
from interface_homology import Core, Loop, Surface, Model
from interface_homology import condense, verify_condensation
from interface_homology.smith import smith, verify_smith, torsion_blocks

core = Core(vertices=1, edges=((0, 0),))
base = Surface(orientable=False, complexity=1, boundary=1,
               circles=(Loop(anchor=0, word=(1, 1)),))
model = Model(core, (base,))
packet = condense(model)
assert verify_condensation(model, packet)
certificate = smith(packet.matrix, cols=packet.relation_count)
assert verify_smith(packet.matrix, certificate,
                    cols=packet.relation_count)
free_rank = packet.rows - len(certificate.diagonal) + packet.free_bulk
torsion = torsion_blocks(certificate.diagonal, packet.two_torsion_bulk)
assert free_rank == 1
assert torsion == ((2, 1),)  # Z plus one Z/2
\end{lstlisting}

The model-to-packet transition authenticates this algebraic input only.
Changing the surrounding code comment to ``a cut of the input knot'' does
not create the missing source evidence.

\section{Counterexamples that the interface must retain}
\label{app:counterexamples}

\subsection{An identical profile can hide arbitrary circle degree}
Let the core be a solid torus. Attach a disc bundle along its annular
vertical boundary, using a primitive slope on the core boundary whose image
in $H_1(K)=\Z$ is $p$ times a generator. This is the familiar solid-torus plus
two-handle description of a punctured lens space when $p\geq2$. Its boundary
is a sphere, but the condensation gives
\[
 C=(p),\qquad H_1=\Z/p.
\]
The base type, number of ports, and Euler characteristic do not determine
$p$. The actual core word does. This example also shows why handle count
alone is not a coefficient-size bound. It is not asserted to arise from
normal cuts with one fixed small source triangulation.

\subsection{Intervals can create homology with no attached circles}
Take two contractible core components and a disc base attached by three
intervals, two to the first component and one to the second. There are three
graph edges, three graph vertices, and one connected component, so
$r_\Gamma=1$. The base and core have no first homology, and $C$ is empty,
but $H_1(X)=\Z$. Omitting the attachment graph would return zero.

\subsection{Genus-zero boundary with several components is not a ball}
The product $S^2\times I\subset S^3$ has two sphere boundary components.
It has $H_1=0$ but $H_2=\Z$ and is not a ball. The boundary shortcut must
retain the number of boundary components, not just their total genus.

\subsection{A homology result is not a marked pattern}
A solid torus and a nontrivial knot exterior both have boundary genus one
and integral homology $(\Z,\Z,0,0)$. Thus the boundary shortcut does not
recognize an unknot from its initial exterior. The further hierarchy and
peripheral/pattern information remain necessary even when every homology
query is cheap.

\end{document}
